% Work/MAT | 1. rocnik | 24. hodina
% Upevneni trojuhelniku, presnych hodnot, kalkulacky a technickych uloh.
% Zdroj: ucebnice Goniometricke funkce.tex a Planimetrie.tex;
% Realisticky.cz 040202, 010120; konkretni zadani a vysledky jsou vlastni.
\documentclass[11pt,a4paper]{article}
% Spolecny styl pro ucitelske pripravy MAT (Work/MAT)
% Verze 1.1, 2026-10-23
\NeedsTeXFormat{LaTeX2e}
\ProvidesPackage{priprava}[2026/10/23 v1.1 Sablona ucitelskych priprav MAT]

% Zakladni typografie a cestina
\RequirePackage[T1]{fontenc}
\RequirePackage[utf8]{inputenc}
\RequirePackage[czech]{babel}
\RequirePackage{lmodern}
\RequirePackage{microtype}
\RequirePackage{amsmath,amssymb}
\RequirePackage{xcolor}
\RequirePackage{enumitem}
\RequirePackage[most]{tcolorbox}
\RequirePackage{tabularx,booktabs,array}
\RequirePackage{fancyhdr}
\RequirePackage{needspace}
\RequirePackage{geometry}
\RequirePackage[hidelinks]{hyperref}

\geometry{a4paper,top=17mm,bottom=18mm,left=23mm,right=23mm}
\definecolor{notegray}{gray}{0.39}
\setlength{\parindent}{0pt}
\setlength{\parskip}{2.5pt plus 1pt minus 0.5pt}
\setlist[itemize]{topsep=2pt,itemsep=1pt,parsep=0pt,leftmargin=1.45em}
\setlist[enumerate]{topsep=2pt,itemsep=2pt,parsep=0pt,leftmargin=1.65em}
\setcounter{secnumdepth}{0}
\renewcommand{\arraystretch}{1.12}

% Zahlavi: vlevo rocnik | globalni tema, vpravo cislo: nazev hodiny
% Zapati zustava zcela prazdne.
\pagestyle{fancy}
\fancyhf{}
\renewcommand{\headrulewidth}{0.35pt}
\renewcommand{\footrulewidth}{0pt}
\setlength{\headheight}{14pt}
\newcommand{\nastavPripravu}[4]{%
  \fancyhead[L]{\small #1\enspace\textbar\enspace #2}%
  \fancyhead[R]{\small #3: #4}%
}

% Minimalisticke cernobile boxy. Kazda definice zvlast.
\tcbset{commonlesson/.style={enhanced,breakable=false,colback=white,colframe=black,
  colbacktitle=white,coltitle=black,boxrule=0.45pt,
  arc=0.8mm,outer arc=0.8mm,left=3mm,right=3mm,top=1.6mm,bottom=1.6mm,
  toptitle=1mm,bottomtitle=0.7mm,fonttitle=\bfseries,
  attach boxed title to top left={xshift=2.5mm,yshift=-1.15mm},
  boxed title style={colback=white,colframe=white,boxrule=0pt,left=1mm,right=1mm}}}
\newtcolorbox{zadaniBox}[1][]{commonlesson,title={Zadání},#1}
\newtcolorbox{prikladBox}[1][]{commonlesson,title={Příklad},#1}
\newtcolorbox{definiceBox}[1][]{commonlesson,title={Definice},#1}

% Spolecny box pro zadani a jeho reseni, bez mezery mezi castmi.
% Pouziti:
%   \begin{prikladReseniBox}
%     Zadani ulohy.
%     \tcblower
%     \begin{align*} ... \end{align*}
%   \end{prikladReseniBox}
% Spodni cast ma jemny sedy podklad a delici linku; lze v ni sazet
% vice radku matematiky. Obe casti maji jen jeden spolecny ramecek.
\tcbset{solutionlesson/.style={
  commonlesson,
  enhanced jigsaw,
  breakable,
  colbacklower=black!3,
  segmentation style={solid,draw=black!25,line width=0.35pt},
  fontlower=\small\color{notegray},
  before lower={\noindent\textbf{Řešení:}\enspace},
  leftlower=3mm,rightlower=3mm
}}
\newtcolorbox{prikladReseniBox}[1][]{solutionlesson,title={Příklad},#1}
\newtcolorbox{zadaniReseniBox}[1][]{solutionlesson,title={Zadání},#1}

% Pomucky pro autora pripravy: sedive poznamky; \reseni pro starsi pripravy
\newcommand{\poznamka}[1]{\par\begingroup\small\color{notegray}\itshape
  \textnormal{\textbf{Poznámka pro učitele.}}\enspace #1\par\endgroup}
\newcommand{\reseni}[1]{\par\begingroup\small\color{notegray}
  \textbf{Očekávané řešení:}\enspace #1\par\endgroup}
\newcommand{\blok}[1]{\needspace{5\baselineskip}\subsection*{#1}}
\newcommand{\vyrok}[1]{\uv{#1}}

\usepackage{tikz}
\nastavPripravu{1. ročník}{Goniometrie}{24}{Procvičení a aplikace}
\begin{document}
\blok{Dnes vše spojíme a téma uzavřeme (0.--5. minuta)}
Nezavádíme nové goniometrické funkce ani novou teorii. Opakujeme:
\emph{přepona, protilehlá a přilehlá odvěsna, Pythagorova věta,
poměry sin/cos/tan/cot, přesné hodnoty a kalkulačka ve stupních}.
\begin{zadaniBox}[title={Postup, který musí zaznít před výpočtem}]
\textbf{1.} Načrtni situaci, vyznač pravý a známý úhel.
\textbf{2.} Urči přeponu a vztah hledané strany ke známému úhlu.
\textbf{3.} Vyber Pythagorovu větu nebo odpovídající goniometrický poměr.
\textbf{4.} Vyjádři neznámou; teprve potom dosaď a spočítej.
\textbf{5.} Ověř jednotky, režim \texttt{DEG} a smysluplnost výsledku.
\end{zadaniBox}

\blok{Výpočet v trojúhelníku (5.--14. minuta)}
\begin{prikladReseniBox}[title={P1: může přepona vyjít menší než odvěsna?},breakable=false]
Odvěsny pravoúhlého trojúhelníku jsou $9\,\mathrm{cm}$ a $12\,\mathrm{cm}$.
Urči přeponu a čtyři goniometrické poměry pro úhel proti kratší odvěsně.
\tcblower
$c=\sqrt{9^2+12^2}=15\,\mathrm{cm}$.
$\sin\alpha=9/15=3/5$, $\cos\alpha=12/15=4/5$,
$\tan\alpha=9/12=3/4$ a $\cot\alpha=12/9=4/3$.
Přepona je nejdelší; poměry nemají jednotku.
\end{prikladReseniBox}
\begin{prikladReseniBox}[title={P2: úhel a délka},breakable=false]
Trojúhelník má $\gamma=90^\circ$, přeponu $c=10\,\mathrm{cm}$
a ostrý úhel $\alpha=30^\circ$. Urči odvěsny bez kalkulačky.
\tcblower
$a=c\sin30^\circ=10\cdot\frac12=5\,\mathrm{cm}$,
$b=c\cos30^\circ=10\cdot\frac{\sqrt3}{2}=5\sqrt3\,\mathrm{cm}$.
Ověření: $5^2+(5\sqrt3)^2=25+75=100=10^2$.
\end{prikladReseniBox}

\blok{Slovní úloha: změř výšku bez výstupu na střechu (14.--24. minuta)}
\begin{prikladReseniBox}[title={P3: výškový úhel a výška očí}]
Pozorovatel stojí na vodorovné rovině $24\,\mathrm m$ od paty svislé budovy.
Oči má ve výšce $1{,}60\,\mathrm m$. Vrchol vidí pod úhlem $32^\circ$
nad vodorovnou rovinou. Urči výšku budovy $H$ na setiny metru.
\tcblower
\begin{center}
\begin{tikzpicture}[scale=.68]
\draw[thick] (-.5,0)--(7.7,0);
\draw[very thick] (6,0)--(6,4.23);
\draw[thick] (0,0)--(0,.48);
\fill (0,.48) circle(2pt);
\draw[dashed] (0,.48)--(6,.48);
\draw[thick] (0,.48)--(6,4.23);
\draw (5.7,.48)--(5.7,.78)--(6,.78);
\draw (.85,.48) arc[start angle=0,end angle=32,radius=.85];
\node at (1.34,.74){$32^\circ$};
\node[below] at (3,0){$24\,\mathrm m$};
\node[left] at (0,.25){\scriptsize $1{,}60\,\mathrm m$};
\node[right] at (6.03,2.4){$H$};
\end{tikzpicture}
\end{center}
Od výšky budovy odečteme úroveň očí:
$\tan32^\circ=\frac{H-1{,}60}{24}$.
Proto $H=1{,}60+24\tan32^\circ\doteq16{,}60\,\mathrm m$.
Přičíst $1{,}60\,\mathrm m$ je nezbytné; jinak by vyšla jen výška nad očima.
\end{prikladReseniBox}

\newpage
\blok{Samostatné procvičení a společná kontrola (24.--41. minuta)}
\begin{zadaniReseniBox}[title={S1: vyber postup sám; a)--c) je jádro}]
\begin{enumerate}[label=\alph*)]
\item Žebřík délky $6\,\mathrm m$ svírá s vodorovnou zemí úhel $70^\circ$.
Jak vysoko dosahuje a jak daleko je jeho pata od stěny?
\item V pravoúhlém trojúhelníku má protilehlá odvěsna k úhlu $\alpha$
délku $1{,}2\,\mathrm m$ a přilehlá $8\,\mathrm m$. Urči $\alpha$.
\item Pomocí přesných hodnot urči $\sin60^\circ\cos30^\circ+
\cos60^\circ\sin30^\circ$. Je výsledek větší než $1$?
\item (Navíc) Je-li $\beta=90^\circ$, $b=20\,\mathrm{cm}$ a $\alpha=40^\circ$,
urči strany $a,c$.
\end{enumerate}
\tcblower
a) Výška $h=6\sin70^\circ\doteq5{,}64\,\mathrm m$,
vodorovná vzdálenost $d=6\cos70^\circ\doteq2{,}05\,\mathrm m$.
b) $\alpha=\arctan\frac{1{,}2}{8}\doteq8{,}53^\circ$.
c) $(\frac{\sqrt3}2)^2+(\frac12)^2=\frac34+\frac14=1$.
d) Přepona je $b$; $a=20\sin40^\circ\doteq12{,}86\,\mathrm{cm}$,
$c=20\cos40^\circ\doteq15{,}32\,\mathrm{cm}$.
\end{zadaniReseniBox}
\poznamka{Základní čas: 5 min opakování, 9 min P1--P2, 10 min P3,
17 min samostatné úlohy a kontrola, 4 min závěrečné ověření.
Úloha S1 d) je dobrovolná a nesmí vytlačit povinnou kontrolu.}

\blok{Závěrečné ověření a vyhlídka na další téma (41.--45. minuta)}
\begin{zadaniBox}[title={Výstupní lístek (bez známky)}]
\textbf{a)} Kdy raději použiji tangens než sinus?
\quad \textbf{b)} Proč musí být u $\arcsin(0{,}5)$ řečeno, jaký úhel hledáme?
\quad \textbf{c)} Co zkontroluji na kalkulačce?
\end{zadaniBox}
\reseni{a) Když znám nebo hledám obě odvěsny vzhledem k úhlu;
b) v této kapitole hledáme ostrý úhel, zatímco goniometrické funkce
mají mimo tuto oblast další řešení; c) režim DEG.}
Přechod k mocninám: druhé mocniny a odmocniny už využíváme v Pythagorově větě;
v příštím bloku jejich význam a početní pravidla prozkoumáme podrobněji.

\blok{Rezerva pro rychlé a domácí práce}
\begin{zadaniReseniBox}[title={R1: technické pokračování},breakable=false]
Při stejné vodorovné vzdálenosti $50\,\mathrm m$ pozorujeme vrchol pod úhlem
$40^\circ$. Oči jsou ve výšce $1{,}50\,\mathrm m$. Urči výšku objektu.
Pak vysvětli, proč není pravda, že \emph{sklon $10\,\%$ znamená úhel $10^\circ$}.
\tcblower
$H=1{,}50+50\tan40^\circ\doteq43{,}45\,\mathrm m$.
Sklon $10\,\%$ znamená poměr převýšení a vodorovné délky $0{,}1$,
tedy $\alpha=\arctan(0{,}1)\doteq5{,}71^\circ$.
\end{zadaniReseniBox}
\begin{zadaniReseniBox}[title={D1: domácí (12--18 minut)}]
\begin{enumerate}[label=\alph*)]
\item Přepona $13\,\mathrm{cm}$, odvěsna $5\,\mathrm{cm}$: dopočti druhou odvěsnu.
\item Pro ostrý úhel $\alpha=45^\circ$ a přeponu $8\,\mathrm{cm}$ urči obě odvěsny přesně.
\item Ve vzdálenosti $15\,\mathrm m$ od paty stožáru naměříme výškový úhel $35^\circ$.
Výšku očí zanedbej a urči výšku stožáru.
\item Který ostrý úhel má $\sin\alpha=0{,}8$? Použij kalkulačku.
\end{enumerate}
\tcblower
a) $\sqrt{13^2-5^2}=12\,\mathrm{cm}$.
b) Obě $8\frac{\sqrt2}{2}=4\sqrt2\,\mathrm{cm}$.
c) $15\tan35^\circ\doteq10{,}50\,\mathrm m$.
d) $\arcsin(0{,}8)\doteq53{,}13^\circ$.
\end{zadaniReseniBox}
\end{document}
